A continuación se presenta un árbol que se ha obtenido al aplicar el algoritmo de \textit{Branch and Bound} de Dakin a un problema de maximizar con dos variables de decisión: $x_1$ y $x_2$.

Cuando se ha obtenido la solución óptima de la relajación lineal de un nodo se indica el valor de la función objetivo y la de la variables de esa solución.


\begin{figure}[h!]
    \begin{center}
    \begin{tikzpicture}[->,>=stealth',level/.style={sibling distance = 5cm/#1,
      level distance = 1.5cm}] 
    \node [arn_n, label=right: \tiny{(3.571, 3.429)}, label=above right: \textbf{\tiny{440}}] {\tiny{$P_{0}$}}
        child{ node [arn_n, label=right: \tiny{(3, 4)}, label=below left: \textbf{\tiny{}}, label=above left: \textbf{\tiny{420}}, label=below left: \textbf{\tiny{}}] {\tiny{$P_{1}$} }
            edge from parent node[above left] {$x_1 \leq 3$}
        	}
        child{ node [arn_n, label=right: \tiny{(4, 2.4)}, label=below right: \textbf{\tiny{}}, label=above right: \textbf{\tiny{428}}] {\tiny{$P_{2}$}} 
            child{ node [arn_n, label=above left:\textbf{\tiny{423.3}}, label=right: \tiny{(4.1, 2)}, label=below right: \textbf{\tiny{}}] {\tiny{$P_{3}$}}
                child{ node [arn_n, label=above right:\textbf{\tiny{410}}, label=right: \tiny{(4, 2)}, label=below right: \textbf{\tiny{}}] {\tiny{$P_{5}$}} edge from parent node[above left] {$x_1 \leq 4$}
                }
                child{ node [arn_n, label=above right:\textbf{\tiny{}}, label=right:  , label=below right: \textbf{\tiny{}}] {\tiny{$P_{6}$}} edge from parent node[above right] {$x_1 \geq 5$}
                }
            edge from parent node[above left] {$x_2 \leq 2$}
                }
            child{ node [arn_n, label= right:\tiny{}, label=above right:\textbf{\tiny{}}, label=below right: \textbf{\tiny{}}] {\tiny{$P_{4}$}} edge from parent node[above right] {$x_2\geq 3$}
                }
        edge from parent node[above right] {$x_1 \geq 4$}
    		}
    ; 
    \end{tikzpicture}
    \end{center}
\end{figure}

Se pide:
\begin{enumerate} 
    \item Indicar la mejor cota superior y la mejor cota inferior tanto del problema propuesto como de cada uno de los subproblemas.
    \item Explicar qué nodos y por qué se pueden podar, si es que se pueden podar.
    \item Indicar cuál sería un siguiente paso aplicando el algoritmo de Dakin.
\end{enumerate}
